\chapter{Glaucoma Surgery}

\section*{What Is This Surgery?}
Glaucoma surgery encompasses a diverse array of ophthalmic procedures primarily designed to achieve and maintain a lower intraocular pressure (IOP) in patients afflicted with glaucoma, a progressive optic neuropathy characterized by irreversible damage to the optic nerve and corresponding visual field loss. These surgical interventions are typically considered when conventional medical therapy with topical eye drops and laser treatments, such as selective laser trabeculoplasty (SLT) or laser peripheral iridotomy (LPI), have proven insufficient to adequately control IOP or prevent the relentless progression of the disease. The fundamental goal across most glaucoma surgeries is to either create a new, alternative drainage pathway for aqueous humor, the fluid that continuously fills and circulates within the anterior chamber of the eye, or to significantly reduce its production by the ciliary body, thereby restoring a healthier and safer pressure balance within the ocular globe. The most common and historically significant incisional procedure is trabeculectomy, which establishes a filtration bleb, while newer techniques include the implantation of glaucoma drainage devices (tube shunts) and a rapidly expanding category of minimally invasive glaucoma surgeries (MIGS). These procedures are critically important in preserving remaining vision and preventing further, irreversible damage to the delicate optic nerve, which is the conduit for visual information to the brain.

\section*{Alternative Names}
\begin{itemize}
  \item Trabeculectomy
  \item Glaucoma filtration surgery
  \item Aqueous shunt implantation
  \item Tube shunt surgery
  \item Minimally Invasive Glaucoma Surgery (MIGS)
  \item Cyclophotocoagulation (CPC)
\end{itemize}

\section*{Relevant Surgical Anatomy}
A thorough understanding of ocular anatomy, particularly the structures involved in aqueous humor dynamics, is paramount for successful glaucoma surgery. Aqueous humor is produced by the ciliary body, located behind the iris in the posterior chamber. It flows through the pupil into the anterior chamber, a space bounded anteriorly by the cornea and posteriorly by the iris and lens. The primary outflow pathway for aqueous humor is through the trabecular meshwork, a sieve-like tissue located in the anterior chamber angle, which drains into Schlemm's canal, a circumferential channel within the limbus. From Schlemm's canal, aqueous humor enters collector channels and then the episcleral venous system. A secondary outflow pathway, the uveoscleral pathway, allows some aqueous to drain through the ciliary muscle into the suprachoroidal space. Glaucoma is often characterized by impaired outflow through the trabecular meshwork, leading to elevated IOP.

Key surgical landmarks and structures include:
\begin{itemize}
  \item \textbf{Limbus:} The junction between the cornea and sclera, a critical reference point for incisions and placement of filtration devices.
  \item \textbf{Conjunctiva:} The transparent membrane covering the sclera, which forms the outer layer of the filtration bleb in trabeculectomy and covers drainage device plates.
  \item \textbf{Sclera:} The tough, white outer coat of the eye, through which a guarded flap is created in trabeculectomy, and to which drainage plates are secured.
  \item \textbf{Trabecular Meshwork \& Schlemm's Canal:} The target for many MIGS procedures, aiming to enhance natural outflow.
  \item \textbf{Ciliary Body:} The site of aqueous humor production, targeted by cyclodestructive procedures.
  \item \textbf{Optic Nerve Head:} The ultimate structure protected by IOP reduction, located posteriorly, receiving its blood supply from the posterior ciliary arteries and central retinal artery.
  \item \textbf{Extraocular Muscles:} Important to consider during tube shunt placement to avoid diplopia.
\end{itemize}
The arterial supply to the anterior segment, including the ciliary body and iris, is primarily from the anterior ciliary arteries and long posterior ciliary arteries, forming the major arterial circle of the iris. Venous drainage is via the vortex veins. While the eye itself lacks conventional lymphatic drainage, the conjunctiva has lymphatic vessels that drain to regional lymph nodes.

\section*{Surgical Principle \& Rationale}
The overarching principle of glaucoma surgery is to achieve a sustained reduction in intraocular pressure (IOP) to a level that is considered safe for the individual patient's optic nerve, thereby halting or significantly slowing the progression of glaucomatous optic neuropathy. This is accomplished through two primary mechanisms: either by enhancing the outflow of aqueous humor from the eye or by decreasing its production. The rationale for surgical intervention arises when medical and laser therapies fail to achieve this target IOP, or when the disease is advanced at presentation, necessitating a more aggressive and definitive pressure-lowering approach.

Specific operative concepts vary by procedure. Trabeculectomy, the gold standard incisional surgery, creates a new, controlled fistula from the anterior chamber to the subconjunctival space. Aqueous humor filters through this opening, collects in a subconjunctival bleb, and is then absorbed by the surrounding tissues, effectively bypassing the compromised natural outflow system. Glaucoma drainage implants (tube shunts) achieve a similar bypass effect by shunting aqueous humor from the anterior chamber through a tube to an equatorial plate secured to the sclera, where a fibrous capsule forms around the plate, regulating outflow. Minimally Invasive Glaucoma Surgeries (MIGS) typically target the natural outflow pathways, either by removing or bypassing portions of the trabecular meshwork, dilating Schlemm's canal, or creating a micro-fistula to the subconjunctival space, aiming for a less invasive approach with a favorable safety profile. Finally, cyclodestructive procedures, such as cyclophotocoagulation, selectively destroy portions of the ciliary body epithelium, which is responsible for aqueous humor production, thereby reducing the total volume of fluid within the eye. The technical goal of all these procedures is to achieve a stable, target IOP that is low enough to halt or slow glaucomatous progression, while minimizing complications such as hypotony (excessively low IOP) or bleb failure due to scarring.

\section*{History \& Surgical Evolution}
The history of glaucoma surgery spans centuries, with early attempts focused on crude methods to relieve the excruciating pain associated with acutely elevated intraocular pressure. The mid-19th century marked a significant milestone when Albrecht von Graefe performed the first iridectomy for acute angle-closure glaucoma, a procedure that involved removing a portion of the iris to improve aqueous flow. However, the modern era of incisional glaucoma surgery truly began in the late 1960s, laying the groundwork for contemporary techniques.

In 1968, J.E. Cairns introduced the modern trabeculectomy, which rapidly became the gold-standard incisional glaucoma operation worldwide. This innovative procedure involved creating a guarded scleral flap and a fistula to allow aqueous humor to filter into the subconjunctival space, forming a filtration bleb. The efficacy and longevity of trabeculectomy were significantly improved in the 1980s and 1990s with the advent of antimetabolites, such as 5-fluorouracil (5-FU) and mitomycin C (MMC), which are applied intraoperatively to modulate wound healing and prevent excessive scarring of the filtration bleb, a common cause of surgical failure. Concurrently, the development of glaucoma drainage implants, pioneered by surgeons like Anthony Molteno in the 1960s and further refined by Ahmed and Baerveldt in the 1990s, provided a crucial alternative for complex cases, particularly those with previous failed surgeries or severe secondary glaucomas. The 21st century has witnessed a rapid expansion in Minimally Invasive Glaucoma Surgery (MIGS) devices and techniques, starting in the early 2000s. These procedures aim for a safer profile with quicker recovery, often targeting specific components of the conventional outflow pathway, and are increasingly utilized earlier in the disease course, sometimes in conjunction with cataract surgery.

\section*{Clinical Indications}
Glaucoma surgery is indicated when medical and laser therapies are insufficient to control intraocular pressure (IOP) or prevent disease progression.

\begin{itemize}
  \item Progressive visual field loss or optic nerve damage despite maximal tolerated medical therapy.
  \item Persistently high intraocular pressure (IOP) that remains above the target range even with multiple glaucoma medications.
  \item Intolerance to glaucoma medications due to significant side effects, or documented patient non-compliance with the prescribed regimen.
  \item Advanced glaucoma at initial presentation, where immediate and significant IOP reduction is critical to prevent imminent vision loss.
  \item Acute angle-closure glaucoma unresponsive to laser peripheral iridotomy and maximal medical management.
  \item Certain types of secondary glaucoma, such as neovascular glaucoma or uveitic glaucoma, which may be refractory to conventional treatments.
  \item Congenital glaucoma, where surgical intervention is often the primary treatment to establish adequate aqueous outflow.
  \item Patients with a high risk of progression, such as those with normal-tension glaucoma showing documented progression despite low IOP.
\end{itemize}

\section*{Clinical Positioning}
Glaucoma surgery is generally considered a secondary or tertiary intervention in the comprehensive management of glaucoma. For the vast majority of patients, the initial therapeutic approach involves a stepwise escalation of medical therapy with topical eye drops to lower intraocular pressure, often followed by laser procedures such as selective laser trabeculoplasty (SLT) for open-angle glaucoma or laser peripheral iridotomy (LPI) for angle-closure glaucoma. Surgical intervention becomes primary only in specific, often severe or congenital, circumstances. For instance, in cases of congenital glaucoma, surgery is typically the first-line treatment to establish adequate aqueous outflow from birth. Similarly, patients presenting with very advanced glaucoma and critically high IOP, or those with certain aggressive secondary glaucomas that are unresponsive to medical management, may proceed directly to surgery as a primary intervention. Minimally Invasive Glaucoma Surgery (MIGS) procedures are increasingly being considered earlier in the disease course, sometimes even concurrently with cataract surgery, blurring the traditional line between primary and secondary intervention for mild to moderate glaucoma. However, for traditional incisional surgeries like trabeculectomy and glaucoma drainage implants, they remain largely reserved for cases where medical and laser therapies have failed to achieve adequate IOP control or halt disease progression, thus serving as a crucial secondary or salvage option.

\section*{Surgical Technique \& Procedure}
Glaucoma surgery is typically performed in a sterile operating room environment, usually under local anesthesia with intravenous sedation, though general anesthesia may be employed for complex cases, pediatric patients, or highly anxious individuals. The patient is positioned supine, and the operative eye is meticulously prepped with antiseptic solutions and draped in a sterile fashion to minimize infection risk.

\begin{itemize}
  \item \textbf{Trabeculectomy:}
    \begin{enumerate}
      \item A conjunctival incision is made, often fornix-based or limbus-based, followed by careful dissection to expose the underlying sclera.
      \item A partial-thickness scleral flap (e.g., 3x3 mm) is meticulously dissected, creating a lamellar pocket.
      \item Antimetabolitic agents, such as mitomycin C (MMC) or 5-fluorouracil (5-FU), are often applied to the scleral bed beneath the flap for a few minutes to inhibit fibroblast proliferation and reduce postoperative scarring, then thoroughly rinsed with balanced salt solution.
      \item A small block of sclera and trabecular meshwork is excised beneath the flap, creating a sclerostomy (a full-thickness opening into the anterior chamber), which serves as the new drainage pathway.
      \item A peripheral iridectomy (removal of a small piece of iris) is performed to prevent iris prolapse into the sclerostomy and subsequent blockage of aqueous outflow.
      \item The scleral flap is carefully reapproximated with fine, often adjustable, sutures, tensioned precisely to control aqueous outflow and prevent immediate hypotony.
      \item The conjunctiva is meticulously closed with fine sutures to create a watertight seal over the filtration bleb, ensuring no leakage.
    \end{enumerate}
  \item \textbf{Glaucoma Drainage Implant (Tube Shunt) Surgery:}
    \begin{enumerate}
      \item A conjunctival incision exposes the sclera, and the implant's plate (e.g., Ahmed, Baerveldt) is secured to the sclera in the equatorial region, typically 8-10 mm posterior to the limbus, often in a superotemporal or inferonasal quadrant.
      \item A small tube is then inserted into the anterior chamber, usually through a clear corneal or limbal incision, or directly through the sclera anterior to the plate.
      \item A patch graft (e.g., donor sclera, pericardium, or synthetic material) is placed over the tube where it crosses the limbus to prevent erosion of the conjunctiva and subsequent exposure.
      \item The conjunctiva is meticulously closed over the plate and tube, ensuring a smooth, watertight surface.
    \end{enumerate}
  \item \textbf{Minimally Invasive Glaucoma Surgery (MIGS):}
    \begin{enumerate}
      \item These procedures are often performed concurrently with cataract surgery through a small clear corneal incision, or as stand-alone procedures.
      \item Depending on the specific device, the surgeon may insert a micro-stent into Schlemm's canal (e.g., iStent, Hydrus Microstent), create a micro-fistula to the subconjunctival space (e.g., Xen Gel Stent), or perform a goniotomy/trabeculotomy to remove or bypass the trabecular meshwork.
      \item The incision is typically small and self-sealing, requiring minimal or no sutures, contributing to a faster recovery.
    \end{enumerate}
\end{itemize}
At the conclusion of any procedure, topical antibiotics and corticosteroids are applied, and a protective eye shield is often placed for patient comfort and protection.

\section*{Preoperative Workup \& Preparation}
Thorough preoperative workup and preparation are essential to optimize surgical outcomes, minimize potential risks, and ensure patient safety during glaucoma surgery. This comprehensive evaluation addresses both ocular and systemic health.

\begin{itemize}
  \item \textbf{Comprehensive Ophthalmic Examination:} This includes a detailed assessment of visual acuity, intraocular pressure (IOP), gonioscopy to evaluate the anterior chamber angle, slit-lamp examination of the anterior segment, dilated fundus examination of the optic nerve, visual field testing, and optical coherence tomography (OCT) of the optic nerve head and retinal nerve fiber layer to document the extent of glaucomatous damage.
  \item \textbf{Systemic Medical Evaluation:} A complete medical history and physical examination are performed to assess the patient's overall health and identify any comorbidities. This often includes blood tests (e.g., complete blood count, electrolytes, renal function, liver function, coagulation profile if indicated), an electrocardiogram (ECG), and chest X-ray, especially for older patients or those with significant systemic comorbidities. Anesthesia clearance is obtained from the patient's primary care physician or an anesthesiologist.
  \item \textbf{Medication Management:}
    \begin{itemize}
      \item Glaucoma medications are typically continued up to the day of surgery, unless specifically instructed otherwise by the surgeon.
      \item Anticoagulants (e.g., warfarin, dabigatran) and antiplatelet agents (e.g., aspirin, clopidogrel) are usually held for a specified period (e.g., 5-7 days for warfarin, 7-10 days for clopidogrel) before surgery to reduce the risk of intraoperative and postoperative bleeding, always after careful consultation with the prescribing physician and considering the patient's cardiovascular risk.
      \item Other systemic medications are reviewed and adjusted as necessary, particularly those affecting blood pressure or blood sugar.
    \end{itemize}
  \item \textbf{NPO Status:} Patients are instructed to be NPO (nil per os - nothing by mouth) for at least 6-8 hours prior to surgery to reduce the risk of aspiration during anesthesia.
  \item \textbf{Prophylactic Antibiotics:} Topical antibiotic drops may be prescribed for 1-3 days before surgery to reduce the bacterial load on the ocular surface and minimize the risk of postoperative infection.
  \item \textbf{Informed Consent:} A detailed discussion with the patient regarding the specific surgical procedure, its anticipated benefits, potential risks and complications, alternative treatment options, and expected postoperative course is conducted, and comprehensive informed consent is obtained.
\end{itemize}

\section*{Contraindications \& Precautions}
Careful consideration of contraindications and precautions is vital to ensure patient safety and optimize the likelihood of surgical success in glaucoma management.

\textbf{Absolute Contraindications:}
\begin{itemize}
  \item Active ocular infection (e.g., bacterial conjunctivitis, keratitis, endophthalmitis) in the operative eye, which must be treated and resolved prior to surgery.
  \item Severe, uncontrolled systemic medical conditions (e.g., unstable angina, recent myocardial infarction, uncontrolled hypertension or diabetes, severe coagulopathy) that pose an unacceptable anesthetic or surgical risk.
  \item Patient refusal or inability to provide informed consent due to cognitive impairment without appropriate legal guardianship.
  \item Known severe allergy to anesthetic agents or medications commonly used during the procedure.
\end{itemize}

\textbf{Relative Contraindications and Precautions:}
\begin{itemize}
  \item \textbf{Extensive Conjunctival Scarring:} Prior ocular surgeries (e.g., retinal detachment repair, multiple previous glaucoma surgeries) or severe inflammatory conditions can lead to significant conjunctival scarring, which may preclude trabeculectomy due to high risk of bleb failure and favor a tube shunt or cyclodestructive procedure.
  \item \textbf{Severe Dry Eye Disease:} Can compromise conjunctival health and bleb integrity, increasing the risk of bleb-related complications such as leakage or infection.
  \item \textbf{Neovascular Glaucoma:} Often associated with significant inflammation and friable new blood vessels, which may increase bleeding risk. Glaucoma drainage implants are often preferred over trabeculectomy in these cases, and panretinal photocoagulation may be required preoperatively.
  \item \textbf{Bleeding Diathesis or Anticoagulant Use:} While often managed by temporarily holding medications, a significant intrinsic bleeding risk may influence surgical choice or require specific intraoperative precautions and meticulous hemostasis.
  \item \textbf{Monocular Patient:} Surgery on a patient's only functional eye carries a higher risk-benefit ratio, necessitating meticulous planning, extensive patient counseling, and often a more conservative surgical approach.
  \item \textbf{Uveitic Glaucoma:} Active intraocular inflammation must be rigorously controlled preoperatively with corticosteroids to minimize postoperative complications and improve surgical success rates.
  \item \textbf{Aphakia/Pseudophakia:} The absence or presence of an intraocular lens may influence the choice of surgical technique and potentially increase the risk of certain complications like hypotony or vitreous prolapse.
  \item \textbf{Young Age:} Pediatric glaucoma surgery requires specialized techniques and carries different risks and considerations compared to adult surgery.
\end{itemize}

\section*{Complications \& Risks}
While generally safe and effective, glaucoma surgery, like any surgical procedure, carries potential risks and complications, which can be broadly categorized into general surgical risks and procedure-specific ocular complications.

\textbf{General Surgical and Anesthesia Risks:}
\begin{itemize}
  \item \textbf{Anesthesia Complications:} These can range from mild (nausea, vomiting, headache) to severe (allergic reactions to medications, respiratory depression, cardiac events such as arrhythmia or myocardial infarction, or stroke).
  \item \textbf{Retrobulbar Hemorrhage:} Bleeding behind the eye during local anesthesia injection, a rare but potentially vision-threatening complication.
  \item \textbf{Pain and Discomfort:} Postoperative pain, though usually mild to moderate and manageable with oral analgesics.
  \item \textbf{Infection:} Although rare, any surgical incision carries a risk of wound infection.
\end{itemize}

\textbf{Ocular and Procedure-Specific Risks:}
\begin{itemize}
  \item \textbf{Bleeding:} Intraoperative or postoperative hemorrhage, including hyphema (blood in the anterior chamber) or vitreous hemorrhage, which can obscure vision.
  \item \textbf{Infection:}
    \begin{itemize}
      \item \textbf{Blebitis:} Infection of the filtration bleb, a serious complication that can rapidly progress to endophthalmitis.
      \item \textbf{Endophthalmitis:} Severe intraocular infection, a rare but devastating complication that can lead to permanent vision loss or even loss of the eye.
    \end{itemize}
  \item \textbf{Hypotony:} Excessively low intraocular pressure, which can lead to a cascade of complications including choroidal effusions, choroidal hemorrhage, hypotony maculopathy (swelling of the macula causing vision loss), or corneal decompensation.
  \item \textbf{Choroidal Effusion/Hemorrhage:} Accumulation of fluid or blood between the choroid and sclera, potentially causing shallowing of the anterior chamber, pain, and significant vision loss.
  \item \textbf{Bleb Leakage:} Leakage of aqueous humor from the filtration bleb, increasing the risk of hypotony and infection.
  \item \textbf{Bleb Failure/Scarring:} The filtration bleb may scar down and cease to function over time, leading to a rise in IOP and requiring further medical, laser, or surgical intervention.
  \item \textbf{Cataract Progression:} Glaucoma surgery, particularly trabeculectomy, can accelerate the development or progression of cataracts, necessitating future cataract surgery.
  \item \textbf{Shallow Anterior Chamber:} Can occur due to hypotony, pupillary block, or choroidal effusions, potentially leading to corneal endothelial damage or peripheral anterior synechiae.
  \item \textbf{Diplopia (Double Vision):} More common with glaucoma drainage implant surgery due to the bulk of the implant plate affecting extraocular muscle function.
  \item \textbf{Vision Loss:} While the primary goal is to preserve vision, severe complications can rarely lead to permanent vision loss.
  \item \textbf{Need for Further Surgery:} Due to inadequate IOP control, bleb failure, or other complications.
\end{itemize}

\section*{Postoperative Recovery Timeline}
Immediately after glaucoma surgery, the patient is transferred to the Post-Anesthesia Care Unit (PACU) for close monitoring of vital signs and recovery from anesthesia. The operative eye is typically covered with a protective eye shield or patch to prevent accidental rubbing or trauma. Pain is usually mild to moderate and effectively managed with oral analgesics. Nausea and vomiting are common post-anesthesia side effects and are treated with antiemetics as needed. Patients are typically discharged home within a few hours, provided they are stable and have received clear instructions.

Over the first 24-48 hours, patients are advised to rest and adhere strictly to a regimen of topical antibiotic and corticosteroid eye drops, often administered frequently (e.g., every 2-4 hours) to prevent infection and control inflammation. Activity restrictions are crucial during this initial period: patients are advised to avoid heavy lifting, bending over, strenuous activities, and rubbing the eye. Frequent follow-up visits, sometimes daily initially, are scheduled to monitor intraocular pressure, assess the surgical site (e.g., bleb appearance for trabeculectomy, tube position for shunts), check for complications like hypotony or shallow anterior chamber, and adjust medications. Over the first 1-2 weeks, the frequency of follow-up visits gradually decreases, and the steroid drop regimen is slowly tapered. Patients are typically advised to wear the eye shield at night and when napping for several weeks. Long-term recovery involves continued regular ophthalmic examinations to monitor IOP, visual fields, and optic nerve health, often for the remainder of their lives, to ensure sustained surgical success and detect any recurrence of glaucoma progression. Diet is usually unrestricted unless specific medical conditions dictate otherwise.

\section*{Surgical Findings \& Pathology}
During glaucoma surgery, the surgeon meticulously documents several key findings. These include the preoperative intraocular pressure (IOP), the status of the anterior chamber angle (e.g., open, closed, presence of synechiae), the appearance of the optic nerve head, and any specific intraoperative challenges or complications encountered. For trabeculectomy, the surgeon notes the size and placement of the scleral flap, the creation of the sclerostomy, the performance of the peripheral iridectomy, and the tension of the scleral flap sutures. For tube shunts, the position of the plate on the sclera and the entry point and patency of the tube into the anterior chamber are critical findings. The immediate postoperative IOP and the initial appearance of the filtration bleb (for trabeculectomy) are also recorded. Regarding pathology, if iris tissue is removed during a peripheral iridectomy, it is typically sent for pathological examination to confirm its identity and rule out any unexpected pathology, though this is rare. Similarly, if a goniotomy or trabeculotomy is performed, the excised trabecular meshwork tissue may be sent for histological analysis. For standard trabeculectomy, the excised scleral block is usually not sent for routine pathology unless there is an unusual appearance or suspicion of a neoplastic process. The pathology report, when applicable, confirms the tissue type and provides microscopic details, but the functional success of the surgery is primarily assessed clinically.

\section*{Expected Postoperative Course}
A normal and successful postoperative course following glaucoma surgery is characterized by a gradual and predictable recovery, leading to stable intraocular pressure (IOP) and preserved visual function. Patients can expect some initial discomfort, including mild pain, foreign body sensation, and redness, which typically resolve within the first few days to weeks and are managed with prescribed eye drops and oral analgesics. Vision may be blurry initially due to inflammation, corneal edema, or changes in IOP, but it generally improves steadily over several weeks. For trabeculectomy patients, the filtration bleb will undergo a maturation process, appearing elevated and diffuse, indicating effective aqueous outflow. The anterior chamber should remain deep and stable. Inflammation, monitored by the surgeon, will gradually subside with the tapering of corticosteroid eye drops. Patients will slowly resume their normal activities, with most restrictions lifted by 4-6 weeks, though heavy lifting and strenuous activities may be restricted longer. The ultimate goal is a stable, target IOP that prevents further optic nerve damage, often allowing for a reduction or cessation of preoperative glaucoma medications, thereby improving the patient's quality of life without significant long-term discomfort or visual compromise.

\section*{Postoperative Care \& Follow-up}
Postoperative care and diligent follow-up are paramount for monitoring surgical success, managing potential complications, and ensuring long-term vision preservation after glaucoma surgery.

\begin{itemize}
  \item \textbf{Frequent Postoperative Examinations:} Initially, visits are very frequent (e.g., daily for the first few days, then weekly, then monthly) to monitor IOP, assess the surgical site (bleb morphology, tube position), check for inflammation, and detect early complications.
  \item \textbf{Suture Lysis/Adjustment:} For trabeculectomy, adjustable sutures may be loosened or cut (suture lysis, often with a laser) in the early postoperative period to enhance aqueous outflow if IOP is too high.
  \item \textbf{Bleb Needling:} If the bleb begins to scar down and IOP rises, a procedure called bleb needling, often with antimetabolite injection, may be performed to re-establish filtration.
  \item \textbf{Topical Medication Adjustment:} The regimen of topical antibiotics and corticosteroids is gradually tapered over several weeks to months, based on the eye's healing and inflammatory response.
  \item \textbf{Long-term IOP Monitoring:} Regular IOP checks are continued lifelong to ensure the pressure remains within the target range.
  \item \textbf{Visual Field Testing and Optic Nerve Imaging:} Periodic visual field tests and optical coherence tomography (OCT) of the optic nerve head and retinal nerve fiber layer are performed to monitor for any signs of glaucoma progression.
  \item \textbf{Activity Resumption Guidance:} Patients receive instructions on when they can safely resume normal activities, including exercise, driving, and work.
  \item \textbf{Warning Signs Education:} Patients are educated on warning signs that require immediate medical attention, such as sudden severe pain, significant vision loss, excessive redness, discharge, or flashes of light.
\end{itemize}
Adherence to the prescribed medication regimen and follow-up schedule is critical for optimal long-term outcomes.

\section*{Epidemiology}
Glaucoma stands as a leading cause of irreversible blindness globally, affecting an estimated 80 million individuals worldwide, with projections indicating this number could exceed 111 million by 2040. Primary open-angle glaucoma (POAG) is the most prevalent form, accounting for approximately 70-80\% of all cases, followed by primary angle-closure glaucoma (PACG). The incidence and prevalence of glaucoma increase significantly with advancing age, particularly after 60 years, and exhibit considerable variation across different ethnic groups; for instance, individuals of African and Hispanic descent show higher rates of POAG, while Asian populations have a higher predisposition to PACG. While medical therapy and laser treatments serve as the initial mainstays of glaucoma management, a substantial proportion of patients, estimated to be between 10-20\%, will eventually require surgical intervention during their lifetime to preserve vision. Trabeculectomy remains the most frequently performed incisional procedure globally, though the adoption of Minimally Invasive Glaucoma Surgery (MIGS) procedures is rapidly increasing, especially in developed countries, due to their favorable safety profile. The overall mortality directly attributable to glaucoma surgery is exceedingly rare; however, significant morbidity, including vision-threatening complications such as endophthalmitis or suprachoroidal hemorrhage, can occur in a small percentage of cases, typically ranging from 1-5\%.

\section*{Outcomes \& Prognosis}
The outcomes and prognosis following glaucoma surgery are generally favorable, with the primary objective being the preservation of existing vision and the prevention of further glaucomatous progression. For trabeculectomy, success rates, typically defined as achieving a target intraocular pressure (IOP) without the need for additional surgery or severe complications, range from 70-90\% at one year, decreasing to 50-70\% at five years, often requiring some adjunctive medical therapy. Glaucoma drainage implants, while often reserved for more complex cases or those with prior surgical failures, demonstrate similar or slightly lower success rates but offer a more predictable outcome in eyes with extensive conjunctival scarring. The Tube Versus Trabeculectomy (TVT) Study, a landmark clinical trial, provided valuable comparative data on the efficacy and safety of these two procedures. MIGS procedures generally offer lower IOP-lowering efficacy compared to trabeculectomy but boast a significantly safer profile and quicker recovery, with success rates varying widely depending on the specific device and careful patient selection. Functional outcomes primarily focus on stable visual acuity and visual fields, which are achieved in the majority of successful cases, indicating that the surgery has effectively halted or significantly slowed the progression of glaucomatous damage. Recurrence of high IOP due to bleb scarring or shunt encapsulation is the most common reason for surgical failure, necessitating further medical or surgical intervention. Prognostic factors influencing success include the type and stage of glaucoma, patient age, prior ocular surgeries, surgeon experience, and patient compliance with postoperative care. While glaucoma surgery does not restore lost vision, it is highly effective in preventing further irreversible vision loss, thereby significantly improving the long-term quality of life for patients.

\section*{References}
\begin{enumerate}
  \item MedlinePlus. Glaucoma surgery. U.S. National Library of Medicine; 2024. Available from: \url{https://medlineplus.gov/glaucomasurgery.html}
  \item American Academy of Ophthalmology. Preferred Practice Pattern: Primary Open-Angle Glaucoma. San Francisco, CA: American Academy of Ophthalmology; 2020.
  \item Weinreb RN, Brandt JD, Garway-Heath DF, et al. Glaucoma. Lancet. 2014;383(9912):190-201.
  \item Schwartz's Principles of Surgery. 11th ed. McGraw-Hill Education; 2019.
\end{enumerate}

\clearpage